\chapter*{Abstract}
\addcontentsline{toc}{chapter}{Abstract (English)}
\markboth{Abstract}{Abstract}

\label{sec:abstract-english}

% \AK{There are some compilation errors pertaining to Hebrew.}

% This thesis\REV{corrected the abstract include path}\REV{clarified that simulation budgets vary by experiment, with 100k bits $\times$ 10 trials applying only to the canonical comparison according to \ref{sec:monte-carlo-ber-estimation})} \REV{updated the MLP training time to ``under 5 seconds'' to match Table~\ref{tbl:table13}.}  presents a comparative study of classical and neural relay strategies for two-hop \REV{"cooperative" removed} communication under a  SISO Rayleigh fading channel model  \REV{\AK{What's a \underline{fixed} Rayleigh fading?} \GZ{"fixed" mean "common" or "under"}}. Nine relay methods are evaluated, including amplify-and-forward (AF), decode-and-forward (DF), multilayer perceptrons (MLPs), generative models (VAE, cGAN), and sequence models (Transformer, Mamba S6, and Mamba-2 SSD). Performance is assessed via Monte Carlo simulation using BER, confidence intervals, and statistical significance testing. On the canonical benchmark, the best neural relays outperform AF at low SNR \REV{but only match} and achieve performance comparable to DF, which remains the strongest \REV{symbol-wise} conventional model-based relaying scheme under the considered assumptions at medium and high SNR while requiring no trainable parameters. \AK{DF is not symbolwise, unless you simple slicing} \REV{"at medium and high SNR regimes with zero trainable parameters" rephrased}. A complexity study shows that performance is driven primarily by model capacity rather than architecture, with a compact 169-parameter MLP matching networks over 100 times larger. Mamba-2 SSD achieves comparable BER to Mamba S6 while training substantially faster. The thesis then examines scenarios beyond the assumptions of classical receivers. For QPSK, the benchmark conclusions remain unchanged, while for 16-QAM a joint 2D classifier removes the BER floor observed with conventional I/Q-split processing. The main contribution is a systematic study of unknown and mismatched channels. In such settings, classical relaying techniques suffer from severe degradation, whereas \REV{"compact" rephrased}a low-complexity learned relay with only 169 trainable parameters learned \AK{What's compact?} \REV{"relaying" removed} approaches the performance of genie-aided Viterbi MLSE without channel knowledge. Across a range of impairments, pilot budgets, and complexity comparisons, the learned relay provides a practical tradeoff between robustness, computational efficiency, and performance. \REV{"The recommended deployment is a hybrid relay that combines neural processing at low SNR with DF at high SNR, a joint-classifier architecture for higher-order modulation, and learned relays for channels with unmodeled memory or distortion" rephrased.} Accordingly, DF is recommended when the assumed channel model closely matches the true channel conditions, whereas learned relays are preferable when channel models are inaccurate, unavailable, or affected by unmodeled memory and distortion. The central conclusion is that learned relaying techniques do not surpass model-aware optimal receivers, \textt{under common settings}, but offer a robust and efficient alternative when accurate channel models are unavailable. 
% \REV{\AK{From a capacity standpoint, DF is optimal. At short blocklengths, DF should be suboptimal.} GZ{clarified}}

% \REV{\AK{In the matched case, according to your claims here, DF is optimal. For mismatched setting, why would one use DF?}\GZ{Agreed. DF is preferred under accurate model matching. Since perfect channel knowledge is rarely available in practice, we also evaluate DF under mismatch to quantify its robustness to modeling errors and compare it with learned relaying.}}

% This thesis presents a systematic comparison of classical and learned relay strategies for two-hop communication over SISO Rayleigh fading channels. Nine methods are evaluated: amplify-and-forward (AF), decode-and-forward (DF), multilayer perceptrons (MLPs), generative models (VAE and cGAN), and sequence models (Transformer, Mamba S6, and Mamba-2 SSD). Performance is quantified through Monte Carlo simulation using bit-error rate (BER), confidence intervals, and statistical significance tests. Under the canonical matched-channel benchmark, the best learned relays outperform AF at low signal-to-noise ratio (SNR) and achieve BER comparable to DF at medium and high SNR. Among the evaluated conventional schemes, DF provides the best BER without trainable parameters. Although DF has strong information-theoretic guarantees in asymptotic capacity settings, these guarantees do not directly establish its optimality under the finite-length, uncoded conditions considered here. The complexity analysis further shows that performance depends primarily on model capacity rather than architectural complexity: a 169-parameter MLP matches networks with more than 100 times as many parameters, while Mamba-2 SSD achieves BER comparable to Mamba S6 with substantially faster training. The study then extends beyond the assumptions of classical model-based receivers. For QPSK, the conclusions of the canonical benchmark remain unchanged. For 16-QAM, a joint two-dimensional classifier eliminates the BER floor observed with conventional independent processing of the in-phase and quadrature components. The principal contribution is a systematic evaluation under unknown and mismatched channels. In these settings, classical model-based methods degrade severely, whereas a learned relay with only 169 trainable parameters approaches the performance of genie-aided Viterbi maximum-likelihood sequence estimation without requiring channel knowledge. These results support a clear deployment strategy: DF is preferred when the assumed channel model accurately represents the true channel, whereas learned relays are advantageous when the channel model is unavailable, inaccurate, or fails to capture memory and distortion. Overall, learned relays do not outperform optimal model-aware receivers under accurately modeled conditions, but provide a robust and computationally efficient alternative under channel uncertainty and mismatch. 
% \AK{Please use AI to polish this abstract.}

This thesis presents a systematic comparison of classical and learned relay strategies for two-hop communication over SISO Rayleigh fading channels. Nine relay methods are considered: amplify-and-forward (AF), decode-and-forward (DF), multilayer perceptrons (MLPs), generative models based on a variational autoencoder (VAE) and a conditional generative adversarial network (cGAN), and sequence models based on the Transformer, Mamba-S6, and Mamba-2-SSD architectures. All methods are assessed through Monte Carlo simulation using bit-error rate (BER), confidence intervals, and statistical significance testing.

Under a canonical matched-channel benchmark, the best learned relays outperform AF at low signal-to-noise ratio (SNR) and attain BER comparable to DF at medium and high SNR, while DF achieves the lowest BER among the conventional schemes without any trainable parameters. Although the literature establishes strong information-theoretic results for DF in asymptotic capacity regimes, these results do not directly imply optimality under the finite-length, uncoded conditions studied here. A complexity analysis indicates that performance is governed primarily by model capacity rather than architectural sophistication: a compact MLP with only 169 trainable parameters attains the BER of networks more than two orders of magnitude larger, and Mamba-2 SSD matches the BER of Mamba S6 at substantially lower training cost.

The robustness of the evaluated methods is then examined in regimes where the modeling assumptions underlying conventional receivers no longer hold. For QPSK, the conclusions of the canonical benchmark are unchanged; for 16-QAM, a joint two-dimensional classifier eliminates the BER floor incurred when the in-phase and quadrature components are processed independently. The principal contribution is a systematic evaluation under unknown and mismatched channels, in which conventional model-based relays degrade substantially on a channel with an unknown intersymbol-interference (ISI) filter, whereas a learned relay with only 169 trainable parameters approaches, without explicit channel knowledge, the performance of genie-aided Viterbi maximum-likelihood sequence estimation (MLSE) matched to that filter.

These results characterize the operating regimes of the evaluated approaches. When the assumed channel model accurately represents the true channel, DF attains the lowest BER among the evaluated per-symbol relays without training \AK{Why? DF cuts the code blocklength by half. Can't this be beaten?}\REV{In the uncoded, per-symbol setting of this thesis, the relay makes a one-shot symbol decision and $\operatorname{sign}(y)$ is the MAP decision, hence optimal among per-symbol relay functions. The blocklength/coding argument does not apply here because no channel code is used; coded block-DF and soft-information (LLR) forwarding, which could indeed do better, are identified as future work (Remark~\ref{rem:df-terminology}) \AK{AF outperforms DF for uncoded symbols for low SNR values. DF is not universally better across all SNR values.}\GZ{Checked directly against this thesis's own data (Table~\ref{tbl:table2}, canonical Rayleigh channel, and its AWGN counterpart): DF's per-symbol slicing has a strictly lower BER than AF at every SNR point evaluated, 0--20~dB, on both channels, so there is no low-SNR crossover within the range studied here. The general point remains valid in the wider literature: \emph{naive} (always-forward) DF can underperform AF at sufficiently low SNR because the relay regenerates and confidently retransmits its own decoding errors, which is why \emph{selective} DF (relaying only above a reliability threshold) is the standard mitigation. That effect simply does not appear within the evaluated SNR range and configuration, so the abstract's claim is accurate as an empirical statement about the results reported here rather than a universal one.}}; under model uncertainty or mismatch, learned relays are more robust to modeling errors. Their principal advantage therefore lies not in surpassing optimal model-aware receivers under accurate modeling, but in preserving performance when the channel model is unavailable or fails to capture channel memory or distortion — for example, when a channel filter introduces intersymbol interference (unmodeled memory), the regime studied in the unknown/mismatched-channel study (Chapter~\ref{sec:unknown-channels}). \AK{MLSE w.r.t.\ what? Did you assume colored noise or a channel filter? Did you use convolutional codes? The System MOdel in Chapter~1.1 appears memoryless with white noise}\REV{MLSE is with respect to the FIR channel filter (3-tap ISI) introduced only in the unknown/mismatched-channel study of Chapter~\ref{sec:unknown-channels}; the noise remains white and no convolutional code is used. The canonical model of Chapter~1 is indeed memoryless — ``channel memory and distortion'' in the abstract refers specifically to that deliberate extension, not the canonical setup. The abstract sentence now names the FIR/ISI channel filter explicitly so that MLSE has a well-defined referent \AK{This si unclear from the abstract. You cannot refer to a model that you've never mentioned}.\GZ{Fixed at the source rather than by cross-reference: the earlier sentence introducing MLSE now names the impairment inline (``a channel with an unknown intersymbol-interference (ISI) filter''), so the abstract is self-contained and MLSE has a well-defined referent without requiring the reader to have already seen Chapter~\ref{sec:unknown-channels}.}}

% \textbf{Keywords:} Cooperative relay communication, deep learning, two-hop relay, Rayleigh fading, MLP, Mamba, Transformer, VAE, conditional GAN, 16-QAM, intersymbol interference, Viterbi MLSE, blind equalization, bit error rate
% 
% This thesis presents a comparative study of classical and neural network-based relay strategies for two-hop cooperative communication, carried out on a single canonical setup fixed in advance: a SISO link on both hops over i.i.d.\ Rayleigh fast fading, complex baseband, BPSK, and uncoded bit-error rate (BER), with the relay function as the only varied axis. Nine relay methods are evaluated: two classical approaches, amplify-and-forward (AF) and decode-and-forward (DF), and seven neural methods spanning supervised learning (MLP, Hybrid SNR-adaptive relay), generative modeling (VAE, conditional GAN with WGAN-GP), and sequence architectures (Transformer, Mamba S6, Mamba-2 SSD). All results use Monte Carlo simulation with 95\% confidence intervals and Wilcoxon significance testing; the canonical relay comparison uses 100{,}000 bits per SNR point over 10 trials, while the channel-validation and unknown-channel studies use their own stated budgets (Section~\ref{sec:monte-carlo-ber-estimation}).

% On the canonical setup, the best neural relays outperform AF at low SNR (0--4 dB) but only match symbol-wise DF, which dominates at medium-to-high SNR with zero parameters. A complexity study reveals an inverted-U relationship between model size and performance: a minimal 169-parameter network matches models 100$\times$ larger, while intermediate sizes overfit; at an equal 3,000-parameter budget, architectural differences largely vanish, indicating that capacity, not architecture, drives performance for this task. Mamba-2 SSD trains substantially faster than Mamba S6 at long context with comparable BER.

% A scoped modulation study shows the canonical findings transfer fully to QPSK via I/Q independence, while 16-QAM exposes a structural BER floor of per-axis I/Q-split processing that a joint 16-class 2D classifier eliminates, matching DF at high SNR. The principal contribution is a systematic study of learned relaying under \emph{unknown and mismatched} channels. On an unknown 3-tap ISI channel the memoryless classical relays are pinned at an analytic BER floor of 0.25, DF worsening as SNR grows, while the same minimal MLP (here 170 parameters) restores reliable relaying within 1--1.5 dB of genie-CSI Viterbi MLSE with no channel-family knowledge. The result is then bounded from every side: it never beats matched Viterbi; on a composite ISI$\times$nonlinearity$\times$phase cascade it tracks $\sim$2 dB behind pilot-aided MLSE with no impairment model; in the posterior-free regime it matches blind CMA while avoiding the instability of decision-directed MLSE; a pilot-budget sweep locates a sharp identifiability crossover governing when the pilot-free learned relay overtakes pilot-aided classical processing; and a complexity analysis shows the learned relay's per-symbol cost is constant in modulation order and channel memory (where Viterbi grows as $M^{L}$) and $30$--$90\times$ faster in wall-clock. The learned relay thus never beats a matched classical receiver, but occupies a well-defined niche of family-agnostic, identification-free, constant-complexity mitigation.

% The recommended deployment is a Hybrid relay combining neural processing at low SNR with DF at high SNR (169 parameters, \textasciitilde0.7 KB, under 5 seconds of training), a joint $N$-class 2D classifier for 16-QAM, and a learned relay wherever the link carries unmodeled memory or bias. The overarching insight is that learned relays never beat the model-aware optimum but are uniquely robust to not knowing which model applies, and that model capacity must be matched to task structure, not merely task size.

% \textbf{Keywords:} Cooperative relay communication, deep learning, two-hop relay, Rayleigh fading, MLP, Mamba, Transformer, VAE, conditional GAN, 16-QAM, intersymbol interference, Viterbi MLSE, blind equalization, bit error rate

% \begin{center}\rule{0.5\linewidth}{0.5pt}\end{center}

% \begin{flushright}
% \begin{otherlanguage}{hebrew}
% \textbf{ארכיטקטורות למידה עמוקה לתקשורת ממסר דו-קפיצתית: מחקר השוואתי של אסטרטגיות ממסר קלאסיות ומבוססות רשתות נוירונים}
% גיל צוקרמן
% חיבור זה הוגש כמילוי חלקי של הדרישות לקבלת תואר מגיסטר למדעים (M.Sc.)
% בית הספר להנדסת חשמל — אוניברסיטת תל אביב
% 2026
% \end{otherlanguage}
% \end{flushright}

\clearpage

%% ── Acknowledgments ──────────────────────────────────────────
\cleardoublepage
\phantomsection
\chapter*{Acknowledgments}
\addcontentsline{toc}{chapter}{Acknowledgments}
\markboth{Acknowledgments}{Acknowledgments}

I would like to thank my supervisors, Dr. Anatoly Khina and Prof. Raja Giryes at Tel Aviv University, for their guidance and support throughout this research.
This work was conducted in the School of Electrical and Electronic Engineering, Tel Aviv University.


